\subsection{Protocolo de revisión académica, industrial y patentaria}
\label{subsec:method-review-protocol}

La brecha que justifica este trabajo se sostiene sobre tres corpus con
protocolos distintos: 59 trabajos académicos canónicos (de 3\,014 descubiertos
por Crossref, OpenAlex y \emph{snowballing}), catorce precedentes industriales
con fuente oficial y 167 registros patentarios (CPC G05D~1/69--1/6987, corte
11-09-2026). Cada corpus documenta algo distinto y ninguno sustituye a los
otros dos: el académico recoge mecanismos con formulación publicada, el
industrial recoge capacidades que el fabricante declara en producto y el
patentario muestra cómo se desplaza la composición temática del corpus de
publicaciones recuperado. Los verbos son deliberados: en esta memoria
«certificar» queda reservado a los predicados formales de los Capítulos~5 y~6,
cuyo dominio y supuestos están enunciados, porque una revisión aporta evidencia
documental y no un certificado. Por la misma razón, el análisis patentario se
enuncia a nivel de publicación y no de familia de patentes ni de invención: no
cuenta invenciones ni mide actividad de protección en sentido estricto. El
protocolo completo, sus límites y la tabla comparada de las tres revisiones se
conservan en el Anexo~\ref{app:review-full-detail}.
